\SourceLeaf{71}
\[
\cos\varphi^r=(1+a)\frac{1}{2^r}.{}^{r}\mathfrak{B}^{(\frac12(r+a))}\cos a\varphi+.\,.\,.+\frac{1}{2^{r-1}}\mathfrak{B}^{(\frac12(r+a)+h)}\cos(2h+a)\varphi+.\,.\,.,
\]

\noindent and in this way two series that at first sight appeared utterly different can be included in a single, quite simple expression.

\Needspace{10\baselineskip}\par\smallskip\hrulefill\smallskip\par

\begin{center}\textbf{REMARKS.}\end{center}

The manuscript of this “Dissertatiuncula” was found in 1893 by Professor \textsc{Wilhelm Meyer} (Göttingen) in the records of the Royal Society of Sciences in Göttingen. E. \textsc{Schering} published it with a series of explanatory remarks in the “Göttinger Nachrichten” for 1893, no.\,15. The original manuscript gives neither the author's name nor the place or date of composition. Comparison of other manuscripts by \textsc{Gauss} with the present manuscript, however, removed any doubt that might otherwise have remained about \textsc{Gauss}'s authorship.

As for the subject matter, \textsc{Gauss} probably studied the works of \textsc{Lagrange} and \textsc{Euler} as early as the time when he attended the Collegium Carolinum in Braunschweig (1792 to 1795). His papers contain various excerpts from works by \textsc{Lagrange} and \textsc{Euler}, probably dating from that time. Among them are, in particular, excerpts from \textsc{Euler}'s Institutiones calculi integralis, part I, chapter 6, on which the Dissertatiuncula directly builds.

Although the contents likewise leave no doubt about \textsc{Gauss}'s authorship, the near-total dependence of the entire development on \textsc{Euler} can be reconciled with this only by supposing that the investigation in question originated either at the very beginning of \textsc{Gauss}'s studies in Göttingen (autumn 1795 to 1798) or even before that period. For \textsc{Euler} solves the problem of expanding $(1+n\cos\varphi)^\nu$ into the series

\[
A + A^{(1)}\cos\varphi + A^{(2)}\cos 2\varphi + A^{(3)}\cos 3\varphi + .\,.\,.
\]

\noindent by giving only $A$ and $A^{(1)}$ directly and referring for the subsequent coefficients to a recurrence formula obtained by a natural device. Dispensing with this device, \textsc{Gauss} obtains the higher coefficients as well, namely $A^{(2)}$,\,.\,. by direct consideration, but otherwise reproduces all of \textsc{Euler}'s initial formulations and subsequent developments.

The citations to \textsc{Kästner} suggest that the final composition of the Dissertatiuncula dates from the beginning of the period of study in Göttingen. The paper cited near the end of no.\,10, by \textsc{Klügel}, belongs to the mathematical Commentationes of the Göttingen Society for 1793 and 94. The Commentationes were indeed published as a collected volume only in 1796, but \textsc{Gauss} probably had access to the individual papers earlier.

The absence of any reference to this paper in the scientific diary, which begins at the end of March 1796, would likewise support dating the Dissertatiuncula to the beginning of 1796 at the latest.

\Needspace{10\baselineskip}
One specific remark is still called for concerning the expansion in no.\,8, $\gamma$. The allegation of an error on \textsc{Euler}'s part is unfounded. \textsc{Gauss} himself made a slip in calculating the equations for $A^{(h)}$, $A^{(h+1)}$ by overlooking the factor $\frac{2}{\BareRoot{(1-nn)}}$ on the right-hand side. In the edition he prepared, \textsc{Schering} added this factor and thereby restored agreement with \textsc{Euler}. The present printing likewise gives the corrected formulas.\hfill\textsc{Fricke}.

\par\smallskip\hrulefill\smallskip\par

